%%% =====================================================================
%%% 5 CÂU TRẮC NGHIỆM (EX) - KHTN 6
%%% Chương 3: Một số vật liệu, nhiên liệu, nguyên liệu, lương thực - thực phẩm
%%% Bài 4: Một số lương thực - thực phẩm thông dụng
%%% =====================================================================

%%%=============EX_1=============%%%
\begin{ex}%[6K3N4-1]
Nhóm chất dinh dưỡng nào sau đây cung cấp năng lượng chủ yếu cho cơ thể và có nhiều trong gạo, ngô, khoai, sắn?
\choice
{Chất đạm (protein)}
{Chất béo (lipid)}
{\True Chất bột đường (carbohydrate)}
{Vitamin và chất khoáng}
\loigiai{
	Các loại lương thực như gạo, ngô, khoai, sắn chứa nhiều tinh bột, thuộc nhóm chất bột đường (carbohydrate). Đây là nguồn cung cấp năng lượng chính cho mọi hoạt động sống của cơ thể. Chất đạm chủ yếu giúp xây dựng và phát triển cơ thể; chất béo là nguồn dự trữ năng lượng; vitamin và chất khoáng giúp điều hòa các hoạt động sống.
}
\end{ex}

%%%=============EX_2=============%%%
\begin{ex}%[6K3H4-2]
Để một thời gian trong điều kiện thường, cơm và bánh mì thường bị mốc, có mùi chua. Hiện tượng này cho thấy lương thực, thực phẩm có tính chất nào sau đây?
\choice
{Không bị biến đổi theo thời gian}
{\True Dễ bị biến đổi, hư hỏng nếu không được bảo quản đúng cách}
{Chỉ hư hỏng khi được đun nóng}
{Càng để lâu thì càng tốt cho sức khỏe}
\loigiai{
	Lương thực, thực phẩm là những chất dễ bị biến đổi (ôi, thiu, mốc, lên men, chua...) dưới tác động của vi khuẩn, nấm mốc, nhiệt độ và độ ẩm trong không khí. Vì vậy nếu không được bảo quản đúng cách, chúng sẽ nhanh hư hỏng và không còn an toàn để sử dụng.
}
\end{ex}

%%%=============EX_3=============%%%
\begin{ex}%[6K3V4-4]
Một gia đình mua nhiều rau, củ, quả tươi nhưng chưa dùng hết trong ngày. Biện pháp bảo quản nào sau đây là hợp lí nhất để giữ rau, củ, quả tươi lâu?
\choice
{Để nơi có ánh nắng mặt trời chiếu trực tiếp}
{\True Cho vào ngăn mát tủ lạnh, bảo quản ở nhiệt độ thấp}
{Ngâm ngập trong nước ấm nhiều giờ}
{Đặt sát bên cạnh bếp lửa cho nhanh khô}
\loigiai{
	Ở nhiệt độ thấp trong ngăn mát tủ lạnh, hoạt động của vi khuẩn và nấm mốc bị làm chậm lại, đồng thời hạn chế quá trình bay hơi nước nên rau, củ, quả giữ được độ tươi lâu hơn. Ánh nắng trực tiếp và nhiệt độ cao gần bếp lửa làm rau củ nhanh héo, hỏng; ngâm nước ấm lâu khiến rau bị úng và mất chất dinh dưỡng.
}
\end{ex}

%%%=============EX_4=============%%%
\begin{ex}%[6K3T4-5]
Khi mua một hộp sữa đóng gói ở siêu thị, thông tin nào sau đây trên nhãn em cần xem trước tiên để bảo đảm an toàn cho sức khỏe?
\choice
{Màu sắc và hình ảnh quảng cáo trên bao bì}
{Giá tiền của sản phẩm}
{\True Hạn sử dụng (ngày sản xuất và ngày hết hạn)}
{Khối lượng tịnh của sản phẩm}
\loigiai{
	Trên nhãn thực phẩm, hạn sử dụng cho biết khoảng thời gian sản phẩm còn an toàn để dùng. Thực phẩm đã quá hạn sử dụng có thể bị biến chất, nhiễm vi khuẩn gây hại, dễ dẫn đến ngộ độc. Vì vậy cần kiểm tra hạn sử dụng trước tiên; giá tiền, màu sắc bao bì hay khối lượng không phản ánh mức độ an toàn của thực phẩm.
}
\end{ex}

%%%=============EX_5=============%%%
\begin{ex}%[6K3C4-3]
Cả nhà bạn An ăn món thịt kho để qua đêm ngoài trời, không cất trong tủ lạnh. Sau bữa ăn vài giờ, mọi người bị đau bụng, buồn nôn và tiêu chảy. Nhận định và cách xử lí nào sau đây là đúng?
\choice
{Đây là hiện tượng bình thường sau khi ăn no, không cần xử lí}
{Nên ăn thêm chính món thịt kho đó để cơ thể quen dần}
{Chỉ cần uống thêm thuốc bổ là sẽ khỏi ngay}
{\True Đây là dấu hiệu ngộ độc thực phẩm; cần cho người bệnh nghỉ ngơi, uống nhiều nước và đưa đến cơ sở y tế nếu triệu chứng nặng}
\loigiai{
	Thức ăn để qua đêm ở nhiệt độ thường (không bảo quản lạnh) là môi trường thuận lợi cho vi khuẩn phát triển và sinh độc tố. Các biểu hiện đau bụng, buồn nôn, tiêu chảy sau khi ăn là dấu hiệu điển hình của ngộ độc thực phẩm. Cách xử lí đúng là: ngừng ăn thức ăn nghi ngờ, cho người bệnh nghỉ ngơi, bù nước (uống nước hoặc dung dịch oresol), theo dõi và nhanh chóng đưa đến cơ sở y tế nếu nôn, tiêu chảy nhiều hoặc mệt lả. Tuyệt đối không ăn tiếp thức ăn đã gây ngộ độc.
}
\end{ex}
